\documentclass[12pt]{article}

\usepackage{amsmath}
\usepackage{amssymb}
\usepackage{graphicx}
\usepackage{hyperref}
\usepackage[margin=1in]{geometry}
\usepackage{natbib}
\usepackage{xcolor}
\usepackage{tcolorbox}
\usepackage{booktabs}
\usepackage{tabularx}
\usepackage{bm}
\tcbuselibrary{breakable,skins}
\bibliographystyle{aasjournal}

\renewcommand{\vec}[1]{\bm{#1}}

\newcommand{\Lag}{\mathcal{L}}
\newcommand{\Ham}{\mathcal{H}}
\newcommand{\pb}[2]{\{#1,#2\}}

\title{Classical Mechanics for Simulation Work:\\
       Lagrange, Legendre, Hamilton, Phase Space,\\
       and Why $N$-body Codes Use Leapfrog}
\author{}
\date{\today}

\begin{document}

\maketitle
\tableofcontents
\newpage

\begin{tcolorbox}[colback=yellow!10!white, colframe=orange!75!black,
                  title={\bfseries Disclaimer}]
These notes were compiled by an AI assistant (Claude, Anthropic) and have
\textbf{not been reviewed by a human domain expert}.  They may contain
errors in derivations, physical reasoning, or numerical estimates.
Cross-check key equations against the primary literature before use.
\end{tcolorbox}
\vspace{0.5em}

%----------------------------------------------------------------------
\section{Purpose and scope}
\label{sec:intro}

Several notes in this collection assume Hamiltonian mechanics without
deriving it.  The kinetic-theory appendix of \texttt{sidm.tex} opens by
writing down Hamilton's equations and immediately builds a flow on
phase space out of them.  The field-theory note \texttt{qft.tex}
performs a Legendre transform on a Lagrangian density in half a page,
on the way to canonical quantisation.  The initial-conditions note
\texttt{cosmo\_ic.tex} mentions that the SWIFT code converts to a
``symplectic integration variable'' at start-up, without saying what
makes a variable symplectic or why one would want that.

This note supplies the missing material.  It is written for someone who
has seen Newton's laws and some vector calculus, and who wants to be
able to read the other notes without stopping.  Three questions
organise it:

\begin{enumerate}
  \item \textbf{Where do Hamilton's equations come from?}  They are not
        a postulate.  They follow from the Lagrangian by a specific
        change of variables, the Legendre transform, and knowing which
        change makes them possible is what stops a reader from
        misapplying them.
  \item \textbf{What is allowed to be called a change of coordinates
        on phase space?}  Not every substitution preserves Hamilton's
        equations.  The ones that do are called canonical, and there
        are three practical ways to test for them.
  \item \textbf{Why does every $N$-body code integrate with leapfrog
        rather than a higher-order Runge--Kutta method?}  The answer is
        a direct consequence of the first two questions, and it is the
        reason this note belongs in a simulation repository.
\end{enumerate}

Sections~\ref{sec:lagrange}--\ref{sec:hamilton} build the formalism,
Sections~\ref{sec:phase_space}--\ref{sec:canonical} develop phase space
and the transformations that respect it, and
Sections~\ref{sec:liouville}--\ref{sec:symplectic} pay off with
Liouville's theorem and symplectic integration.

Throughout, $q$ denotes a generalised coordinate and $p$ its conjugate
momentum.  For a system with $n$ degrees of freedom there are $n$ of
each, written $q_i$ and $p_i$ with $i = 1,\ldots,n$; a single particle
moving in three dimensions has $n = 3$, and $N$ particles in three
dimensions have $n = 3N$.  Sums over $i$ are written explicitly.

%----------------------------------------------------------------------
\section{Lagrangian mechanics in brief}
\label{sec:lagrange}

This section is a summary, included so that the note stands alone.  The
same derivation is carried out in more detail, and then generalised to
fields, in \texttt{qft.tex}~\S3.

\paragraph{Generalised coordinates.}  Newton's second law is written in
Cartesian components, which becomes awkward as soon as the system has
constraints.  A bead on a wire has three Cartesian coordinates but only
one genuine degree of freedom, and enforcing the constraint requires
tracking the force the wire exerts, which is of no physical interest.
The Lagrangian formulation avoids this by using \emph{generalised
coordinates}: any set of $n$ numbers $q_1,\ldots,q_n$ that specifies
the configuration of the system completely, with $n$ equal to the
number of true degrees of freedom.  For the bead, $n = 1$ and $q$ is
the distance along the wire.  No constraint forces appear.

\paragraph{The action and its stationary point.}  Given a Lagrangian
$L(q,\dot q,t)$, the action of a candidate trajectory $q(t)$ between
two fixed times is
%
\begin{equation}
  S[q] = \int_{t_1}^{t_2} L\bigl(q(t),\dot q(t),t\bigr)\,dt .
  \label{eq:action}
\end{equation}
%
The physical trajectory is the one for which $S$ is stationary against
small variations $q(t) \to q(t) + \delta q(t)$ that vanish at the two
endpoints.  Varying \eqref{eq:action}, integrating the $\delta\dot q$
term by parts, and using $\delta q(t_1) = \delta q(t_2) = 0$ to kill
the boundary term gives
%
\begin{equation}
  \delta S = \int_{t_1}^{t_2}\!
    \left[\frac{\partial L}{\partial q}
        - \frac{d}{dt}\frac{\partial L}{\partial\dot q}\right]
    \delta q \; dt .
\end{equation}
%
Requiring this to vanish for every choice of $\delta q(t)$ forces the
bracket to vanish at each instant, which is the \textbf{Euler--Lagrange
equation}
%
\begin{equation}
  \frac{d}{dt}\frac{\partial L}{\partial\dot q_i}
    - \frac{\partial L}{\partial q_i} = 0 ,
  \qquad i = 1,\ldots,n .
  \label{eq:euler_lagrange}
\end{equation}
%
For a particle in a potential, $L = \tfrac12 m\dot q^2 - V(q)$, this
returns $m\ddot q = -V'(q)$, which is Newton's second law.  The
Lagrangian is therefore not new physics; it is a repackaging that makes
symmetries and coordinate changes easier to handle.

\paragraph{Where this stops being enough.}  Equation
\eqref{eq:euler_lagrange} is a set of $n$ second-order differential
equations in the variables $q_i$.  That is a perfectly good description
of the motion, and for solving a specific problem it is often the
easiest one.  It is a poor starting point for three things we want
later: describing a statistical ensemble of systems, changing variables
freely, and constructing numerical integrators with good long-term
behaviour.  All three want $2n$ first-order equations in a space where
position and momentum appear on the same footing.  Producing that is
the job of the next two sections.

%----------------------------------------------------------------------
\section{The Legendre transform and the Hamiltonian}
\label{sec:legendre}

\paragraph{The problem to solve.}  We want to trade the variable
$\dot q$ for something else, so that the description uses $(q,p)$
instead of $(q,\dot q)$.  The naive approach --- define $p$ somehow and
substitute --- runs into a difficulty that is worth seeing before the
solution is presented.

A function $L(q,\dot q)$ has the differential
%
\begin{equation}
  dL = \frac{\partial L}{\partial q}\,dq
     + \frac{\partial L}{\partial\dot q}\,d\dot q .
  \label{eq:dL}
\end{equation}
%
The presence of $d\dot q$ on the right is what tells us that $L$ is
naturally a function of $\dot q$.  If we want a function whose
differential contains $dp$ instead, we must build a new function whose
$d\dot q$ terms cancel.  The construction that does this is the
\emph{Legendre transform}, and it is a general tool: the same move
turns internal energy into enthalpy in thermodynamics, by trading
volume for pressure.

\paragraph{The canonical momentum.}  Define
%
\begin{equation}
  p_i \equiv \frac{\partial L}{\partial\dot q_i} ,
  \label{eq:canonical_momentum}
\end{equation}
%
called the \textbf{canonical momentum} conjugate to $q_i$.  For
$L = \tfrac12 m\dot q^2 - V(q)$ this gives $p = m\dot q$, the familiar
momentum, but the definition is more general and the two need not
agree.  A charged particle in a magnetic field has
$p = m\dot q + qA$, and a particle described in spherical coordinates
has an angular canonical momentum with units of angular momentum, not
of momentum.  It is the quantity defined by
\eqref{eq:canonical_momentum} that belongs in what follows, not the
Newtonian $m\dot q$.

\paragraph{The transform.}  Define the \textbf{Hamiltonian}
%
\begin{equation}
  \Ham(q,p,t) \equiv \sum_i p_i\dot q_i - L(q,\dot q,t) ,
  \label{eq:hamiltonian_def}
\end{equation}
%
where every remaining $\dot q_i$ on the right is to be re-expressed in
terms of $q$ and $p$ by inverting \eqref{eq:canonical_momentum}.  That
inversion is possible whenever $\partial^2 L/\partial\dot q_i\partial\dot q_j$
is an invertible matrix, which holds for every system in these notes.

Now take the differential of \eqref{eq:hamiltonian_def} and substitute
\eqref{eq:dL}:
%
\begin{equation}
  d\Ham = \sum_i\Bigl(
     \underbrace{p_i\,d\dot q_i}_{\text{from } p\dot q}
   + \dot q_i\,dp_i
   - \frac{\partial L}{\partial q_i}\,dq_i
   - \underbrace{\frac{\partial L}{\partial\dot q_i}\,d\dot q_i}_{= \; p_i\,d\dot q_i}
     \Bigr)
   - \frac{\partial L}{\partial t}\,dt .
  \label{eq:dH_cancellation}
\end{equation}
%
The two marked terms are equal by the definition
\eqref{eq:canonical_momentum} of $p_i$, and they carry opposite signs,
so they cancel.  This cancellation is the entire point of the
construction.  What remains is
%
\begin{equation}
  d\Ham = \sum_i\Bigl(\dot q_i\,dp_i
        - \frac{\partial L}{\partial q_i}\,dq_i\Bigr)
        - \frac{\partial L}{\partial t}\,dt ,
  \label{eq:dH}
\end{equation}
%
which contains $dq$, $dp$ and $dt$ but no $d\dot q$.  The Hamiltonian
is therefore a function of $q$, $p$ and $t$, exactly as we wanted.

\paragraph{Worked example.}  For the one-dimensional harmonic
oscillator, $L = \tfrac12 m\dot q^2 - \tfrac12 m\omega^2 q^2$.  Then
$p = \partial L/\partial\dot q = m\dot q$, so $\dot q = p/m$, and
%
\begin{equation}
  \Ham = p\dot q - L
       = \frac{p^2}{m} - \left(\frac{p^2}{2m}
         - \tfrac12 m\omega^2 q^2\right)
       = \frac{p^2}{2m} + \tfrac12 m\omega^2 q^2 ,
  \label{eq:sho_hamiltonian}
\end{equation}
%
which is the total energy.  That is the usual situation but not a
theorem: $\Ham$ equals the energy whenever the coordinates do not
depend on time explicitly and the potential does not depend on
velocity.  In a rotating frame, for instance, the two differ.

%----------------------------------------------------------------------
\section{Hamilton's equations}
\label{sec:hamilton}

Equation \eqref{eq:dH} was obtained by manipulating $L$.  We may also
write the differential of $\Ham$ directly from the fact that it is a
function of $q$, $p$ and $t$:
%
\begin{equation}
  d\Ham = \sum_i\Bigl(\frac{\partial\Ham}{\partial q_i}\,dq_i
        + \frac{\partial\Ham}{\partial p_i}\,dp_i\Bigr)
        + \frac{\partial\Ham}{\partial t}\,dt .
  \label{eq:dH_direct}
\end{equation}
%
Two expressions for the same differential must agree coefficient by
coefficient, since $dq_i$, $dp_i$ and $dt$ are independent.  Comparing
\eqref{eq:dH} with \eqref{eq:dH_direct} gives
%
\begin{equation}
  \frac{\partial\Ham}{\partial p_i} = \dot q_i ,
  \qquad
  \frac{\partial\Ham}{\partial q_i} = -\frac{\partial L}{\partial q_i} ,
  \qquad
  \frac{\partial\Ham}{\partial t} = -\frac{\partial L}{\partial t} .
\end{equation}
%
The middle relation is not yet in a useful form, but the
Euler--Lagrange equation \eqref{eq:euler_lagrange} says
$\partial L/\partial q_i = d(\partial L/\partial\dot q_i)/dt = \dot p_i$.
Substituting gives \textbf{Hamilton's equations}:
%
\begin{equation}
  \boxed{\;
  \dot q_i = \frac{\partial\Ham}{\partial p_i} ,
  \qquad
  \dot p_i = -\frac{\partial\Ham}{\partial q_i} .
  \;}
  \label{eq:hamilton}
\end{equation}
%
These are $2n$ first-order equations, replacing the $n$ second-order
equations of \eqref{eq:euler_lagrange}.  The count of independent
initial conditions is the same, $2n$, as it must be.

\paragraph{Reading the two equations.}  The first says that the
derivative of $\Ham$ with respect to a momentum gives the corresponding
velocity, which for $\Ham = p^2/2m + V$ returns $\dot q = p/m$.  The
second says that the derivative with respect to a coordinate gives
minus the rate of change of momentum, which returns
$\dot p = -V'(q)$.  Together they are Newton's law, split into two
halves and made symmetric.

\paragraph{The minus sign.}  The asymmetry between the two equations
--- one has a minus sign, the other does not --- is not cosmetic and
will not go away.  It is responsible for the cancellation that proves
Liouville's theorem in \S\ref{sec:liouville}, and it is the reason
phase space carries a structure that ordinary space does not.

\paragraph{Example: the harmonic oscillator.}  With
\eqref{eq:sho_hamiltonian},
%
\begin{equation}
  \dot q = \frac{\partial\Ham}{\partial p} = \frac{p}{m} ,
  \qquad
  \dot p = -\frac{\partial\Ham}{\partial q} = -m\omega^2 q .
\end{equation}
%
Differentiating the first and substituting the second gives
$\ddot q = -\omega^2 q$, the equation of simple harmonic motion at
angular frequency $\omega$, independently of the mass.  We will return
to this system repeatedly, because every feature of the formalism is
visible in it.

\paragraph{Conservation of the Hamiltonian.}  If $\Ham$ has no explicit
time dependence then it is conserved along the motion.  The proof is
one line:
%
\begin{equation}
  \frac{d\Ham}{dt}
   = \sum_i\Bigl(\frac{\partial\Ham}{\partial q_i}\dot q_i
   + \frac{\partial\Ham}{\partial p_i}\dot p_i\Bigr)
   = \sum_i\Bigl(\frac{\partial\Ham}{\partial q_i}
       \frac{\partial\Ham}{\partial p_i}
   - \frac{\partial\Ham}{\partial p_i}
       \frac{\partial\Ham}{\partial q_i}\Bigr) = 0 ,
  \label{eq:H_conserved}
\end{equation}
%
where Hamilton's equations were used in the second step and the two
terms cancel because of the minus sign just discussed.  When $\Ham$ is
the energy, this is conservation of energy.

%----------------------------------------------------------------------
\section{Phase space}
\label{sec:phase_space}

\paragraph{The space.}  Collect the coordinates and momenta into one
list,
%
\begin{equation}
  \vec\Gamma = (q_1,\ldots,q_n,p_1,\ldots,p_n) ,
  \label{eq:gamma}
\end{equation}
%
a point in a $2n$-dimensional space called \textbf{phase space}.  A
single point of this space specifies the state of the system
completely: give me $\vec\Gamma$ at one instant and Hamilton's
equations determine the entire past and future.  That is not true of
configuration space, where knowing all the positions is not enough.

\paragraph{Two names that are easy to confuse.}  For a system of $N$
particles in three dimensions, phase space has $6N$ dimensions and is
conventionally called \textbf{$\Gamma$-space}.  The $6$-dimensional
space belonging to one particle on its own is called
\textbf{$\mu$-space}.  The letters come from the Ehrenfests' 1911
review article: $\Gamma$ for \emph{Gas}, the whole system, and $\mu$
for \emph{Molek\"ul}, a single molecule.  Knowing the origin of the
letters makes the pair difficult to mix up.  A single point of
$\Gamma$-space corresponds to $N$ points of $\mu$-space, one per
particle.

\paragraph{The flow.}  Read Hamilton's equations \eqref{eq:hamilton}
not as a statement about one trajectory but as a statement about every
point of phase space at once.  At each $\vec\Gamma$ the right-hand
sides are definite numbers, and together they form a
$2n$-component vector
%
\begin{equation}
  \dot{\vec\Gamma} = (\dot q_1,\ldots,\dot q_n,
                      \dot p_1,\ldots,\dot p_n)
    = \Bigl(\frac{\partial\Ham}{\partial p_1},\ldots,
            -\frac{\partial\Ham}{\partial q_n}\Bigr) .
  \label{eq:gamma_dot}
\end{equation}
%
Assigning a velocity to every point of a space is called a
\emph{velocity field}, and the motion it generates is called a
\textbf{flow}.  The picture to hold is a fluid filling phase space and
streaming along, with each fluid element tracing out one solution of
the equations of motion.

Two properties of this flow are used constantly.  First, the velocity
at a point depends only on the location of that point, not on the
history of how the system arrived there.  Second, it follows that a
trajectory is fixed completely by its starting point, so two distinct
trajectories can never pass through the same point of phase space:
that point would have to carry two different velocities at once.
Trajectories in phase space neither cross nor merge.

\paragraph{Phase portrait of the oscillator.}  For the harmonic
oscillator, $\Ham = p^2/2m + \tfrac12 m\omega^2 q^2$ is conserved, so
each trajectory stays on a curve of constant $\Ham$.  Those curves are
ellipses in the $(q,p)$ plane, with semi-axes
$\sqrt{2\Ham/m\omega^2}$ along $q$ and $\sqrt{2m\Ham}$ along $p$.  The
motion runs clockwise: at the rightmost point $q > 0$, $p = 0$, and the
second of Hamilton's equations gives $\dot p < 0$.  Every trajectory is
closed, which is the phase-space statement that the motion is periodic.

%----------------------------------------------------------------------
\section{Poisson brackets}
\label{sec:poisson}

\paragraph{Definition.}  For any two functions $A(q,p,t)$ and
$B(q,p,t)$ on phase space, the \textbf{Poisson bracket} is
%
\begin{equation}
  \pb{A}{B} \equiv \sum_{i=1}^{n}\left(
     \frac{\partial A}{\partial q_i}\frac{\partial B}{\partial p_i}
   - \frac{\partial A}{\partial p_i}\frac{\partial B}{\partial q_i}
  \right) .
  \label{eq:poisson_def}
\end{equation}
%
Different books order the two terms differently, which flips the
overall sign; the convention above is the one used in
\texttt{sidm.tex}, and it is fixed by requiring
$\pb{q}{p} = +1$.  When comparing with another source, check that
bracket first.

\paragraph{Properties.}  Directly from \eqref{eq:poisson_def}:
%
\begin{align}
  \pb{A}{B} &= -\pb{B}{A}
    &&\text{(antisymmetry, so } \pb{A}{A} = 0),
  \label{eq:pb_antisym}\\
  \pb{A}{\alpha B + \beta C} &= \alpha\pb{A}{B} + \beta\pb{A}{C}
    &&\text{(linearity, for constants } \alpha,\beta),\\
  \pb{A}{BC} &= \pb{A}{B}C + B\pb{A}{C}
    &&\text{(Leibniz rule)},\\
  0 &= \pb{A}{\pb{B}{C}} + \pb{B}{\pb{C}{A}} + \pb{C}{\pb{A}{B}}
    &&\text{(Jacobi identity)} .
\end{align}
%
The Jacobi identity is tedious to verify by hand and is stated here
only because it is what makes the brackets an algebraic structure
rather than a notational convenience.

\paragraph{The fundamental brackets.}  Applying
\eqref{eq:poisson_def} to the coordinates and momenta themselves gives
%
\begin{equation}
  \pb{q_i}{q_j} = 0 ,
  \qquad
  \pb{p_i}{p_j} = 0 ,
  \qquad
  \pb{q_i}{p_j} = \delta_{ij} .
  \label{eq:fundamental_brackets}
\end{equation}
%
Check the last one: only the term with the matching index survives in
the sum, and it contributes
$(\partial q_i/\partial q_i)(\partial p_j/\partial p_j) = 1$ when
$i = j$, with the second term vanishing because $q$ does not depend on
$p$.  These three relations are the algebraic content of the statement
``$q$ and $p$ are conjugate,'' and \S\ref{sec:canonical} will use them
as the definition of a good coordinate system.

\paragraph{Time evolution in one formula.}  Let $f(q,p,t)$ be any
quantity built from the coordinates and momenta.  Its rate of change
along the motion is, by the chain rule and then Hamilton's equations,
%
\begin{equation}
  \frac{df}{dt}
   = \frac{\partial f}{\partial t}
   + \sum_i\Bigl(\frac{\partial f}{\partial q_i}\dot q_i
   + \frac{\partial f}{\partial p_i}\dot p_i\Bigr)
   = \frac{\partial f}{\partial t} + \pb{f}{\Ham} .
  \label{eq:evolution}
\end{equation}
%
This single equation contains all of the dynamics.  Setting $f = q_i$
recovers the first of Hamilton's equations, and $f = p_i$ recovers the
second; setting $f = \Ham$ recovers the conservation law
\eqref{eq:H_conserved}, since $\pb{\Ham}{\Ham} = 0$ by antisymmetry.

\paragraph{Conserved quantities.}  Equation \eqref{eq:evolution} gives
an immediate test: a quantity with no explicit time dependence is
conserved if and only if $\pb{f}{\Ham} = 0$.  Finding constants of the
motion is thereby reduced to computing brackets, which is mechanical,
rather than to inspecting the solution, which is not.

%----------------------------------------------------------------------
\section{Canonical transformations}
\label{sec:canonical}

\paragraph{The question.}  Suppose we replace $(q,p)$ by new variables
$(Q,P)$, each a known function of the old ones.  Hamilton's equations
were derived in the original variables.  Do they still hold in the new
ones?

Not in general.  This is the single most common source of error in
elementary Hamiltonian mechanics, because the equations look like
ordinary partial derivatives that should transform like anything else.
A transformation for which \eqref{eq:hamilton} does survive is called
\textbf{canonical}, and there are three practical tests for it.

\paragraph{Test 1: the fundamental Poisson brackets.}  A
transformation is canonical if and only if the new variables satisfy
the same relations \eqref{eq:fundamental_brackets} as the old ones,
%
\begin{equation}
  \pb{Q_i}{Q_j}_{q,p} = 0 ,
  \qquad
  \pb{P_i}{P_j}_{q,p} = 0 ,
  \qquad
  \pb{Q_i}{P_j}_{q,p} = \delta_{ij} ,
  \label{eq:canonical_test_pb}
\end{equation}
%
where the subscript is a reminder that the derivatives inside the
brackets are taken with respect to the \emph{old} variables.  This is
the criterion to remember.  It is short to check, it needs no
machinery beyond \eqref{eq:poisson_def}, and unlike the next test it
remains correct for any number of degrees of freedom.

\paragraph{Test 2: the Jacobian, with a caveat.}  Assemble the matrix
of first derivatives of the new variables with respect to the old,
%
\begin{equation}
  M = \frac{\partial(Q,P)}{\partial(q,p)},
  \qquad
  M_{ab} = \frac{\partial(Q,P)_a}{\partial(q,p)_b} ,
  \label{eq:jacobian_matrix}
\end{equation}
%
a $2n\times 2n$ matrix.  For \emph{one} degree of freedom, the
transformation is canonical exactly when $\det M = 1$, which says the
map preserves areas in the $(q,p)$ plane.  This is easy to compute and
easy to picture, which is why it is often quoted.

For more than one degree of freedom it is \textbf{necessary but not
sufficient}.  A counterexample with $n = 2$ makes this concrete.  Take
%
\begin{equation}
  Q_1 = \lambda q_1, \quad P_1 = \lambda p_1, \quad
  Q_2 = q_2/\lambda, \quad P_2 = p_2/\lambda .
\end{equation}
%
The Jacobian matrix is diagonal with entries
$\lambda,1/\lambda,\lambda,1/\lambda$, so $\det M = 1$ and the
four-dimensional volume is preserved.  But
$\pb{Q_1}{P_1} = \lambda^2 \neq 1$ unless $\lambda = \pm 1$, so the
transformation is not canonical.  What has happened is that the first
pair's area was stretched by $\lambda^2$ and the second pair's shrunk
by the same factor, leaving the total volume alone.  Volume
preservation is blind to that trade; the brackets are not.

The correct matrix statement, which does hold for any $n$, is the
\emph{symplectic condition} $M \Omega M^{\mathsf T} = \Omega$, where
$\Omega$ is the $2n\times2n$ matrix with $+I$ in the upper-right
block, $-I$ in the lower-left, and zeros elsewhere.  Written out, it
is equivalent to \eqref{eq:canonical_test_pb}.

\paragraph{Test 3: build it from a generating function.}  The third
option avoids testing altogether, by constructing transformations that
cannot fail.  Pick any function $F_2(q,P,t)$ of the old coordinates and
the new momenta, and define the transformation implicitly by
%
\begin{equation}
  p_i = \frac{\partial F_2}{\partial q_i} ,
  \qquad
  Q_i = \frac{\partial F_2}{\partial P_i} ,
  \qquad
  K = \Ham + \frac{\partial F_2}{\partial t} ,
  \label{eq:generating_F2}
\end{equation}
%
where $K$ is the Hamiltonian in the new variables.  Every $F_2$
whatsoever produces a canonical transformation.  Three other versions
exist, $F_1(q,Q)$, $F_3(p,Q)$ and $F_4(p,P)$, differing in which pair
of old and new variables is treated as independent; $F_2$ is the one
needed most often because it includes the identity transformation
$F_2 = \sum_i q_i P_i$.

\paragraph{Worked example: rescaling the oscillator.}  The elliptical
phase portrait of \S\ref{sec:phase_space} is awkward because its axes
depend on $m$ and $\omega$.  Consider
%
\begin{equation}
  Q = \sqrt{m\omega}\,q ,
  \qquad
  P = p/\!\sqrt{m\omega} .
  \label{eq:sho_rescale}
\end{equation}
%
All three tests agree that this is canonical.

\emph{By brackets:}
$\pb{Q}{P} = (\partial Q/\partial q)(\partial P/\partial p)
- (\partial Q/\partial p)(\partial P/\partial q)
= \sqrt{m\omega}\cdot(m\omega)^{-1/2} - 0 = 1$.

\emph{By Jacobian:} the matrix is
$\mathrm{diag}\bigl(\sqrt{m\omega},(m\omega)^{-1/2}\bigr)$, whose
determinant is $1$.  With $n = 1$ this suffices.

\emph{By generating function:} take
$F_2(q,P) = \sqrt{m\omega}\,qP$.  Then
$Q = \partial F_2/\partial P = \sqrt{m\omega}\,q$ and
$p = \partial F_2/\partial q = \sqrt{m\omega}\,P$, which inverts to
$P = p/\!\sqrt{m\omega}$, reproducing \eqref{eq:sho_rescale} with no
verification needed.

In the new variables the Hamiltonian \eqref{eq:sho_hamiltonian}
becomes
%
\begin{equation}
  \Ham = \tfrac12\,\omega\,(Q^2 + P^2) ,
\end{equation}
%
whose level sets are circles, and Hamilton's equations read
$\dot Q = \omega P$, $\dot P = -\omega Q$: a rigid rotation of the
$(Q,P)$ plane at angular velocity $\omega$.  Note also that $Q$ and $P$
now carry the same physical dimensions, so the phase plane has a
genuine geometry, and the product $QP$ has the dimensions of action.

\paragraph{The trap.}  It is tempting to simplify by rescaling the
momentum alone, $Q = q$ and $P = p/m\omega$, since that also turns the
ellipse into a circle.  It is not canonical: its Jacobian is
$1/m\omega$, and $\pb{Q}{P} = 1/m\omega \neq 1$.  Applying
\eqref{eq:hamilton} anyway would give
$\dot q = \partial\Ham/\partial P = m\omega^2 P$, which is wrong by the
factor $m\omega$ and is dimensionally a force rather than a velocity.
The general rule for a pure rescaling $Q = \lambda q$, $P = \mu p$ is
that the two factors must be reciprocals, $\lambda\mu = 1$.

%----------------------------------------------------------------------
\section{Liouville's theorem}
\label{sec:liouville}

\paragraph{Statement.}  The flow \eqref{eq:gamma_dot} generated by
Hamilton's equations preserves phase-space volume.  Take any region
$R_0$ of phase space, let every point of it move along the flow, and
call the region occupied at time $t$ by the moving points $R_t$.  Then
$R_t$ has the same $2n$-dimensional volume as $R_0$, for every $t$.
The shape may distort without limit --- a compact blob is typically
drawn out into a thin filament wound many times around the accessible
region --- but the volume enclosed does not change.

\paragraph{Proof.}  A flow changes volumes at a rate given by the
divergence of its velocity field, exactly as in ordinary fluid
mechanics.  Compute that divergence, grouping the $2n$ terms into
coordinate and momentum pairs and substituting Hamilton's equations:
%
\begin{equation}
  \nabla_{\!\vec\Gamma}\!\cdot\!\dot{\vec\Gamma}
   = \sum_{i=1}^{n}\left(
       \frac{\partial\dot q_i}{\partial q_i}
     + \frac{\partial\dot p_i}{\partial p_i}\right)
   = \sum_{i=1}^{n}\left(
       \frac{\partial^2\Ham}{\partial q_i\,\partial p_i}
     - \frac{\partial^2\Ham}{\partial p_i\,\partial q_i}\right)
   = 0 .
  \label{eq:liouville_theorem}
\end{equation}
%
The two second derivatives are equal because the order of partial
differentiation does not matter, and they carry opposite signs because
of the minus sign in the second of Hamilton's equations.  The
cancellation happens separately for each $i$.  A divergence-free
velocity field is incompressible, which is the statement of the
theorem.

Notice what the proof used: Hamilton's equations and the equality of
mixed partial derivatives, and nothing else.  No probability, no
statistics, and no distribution function appear.  Liouville's theorem
is a geometric property of Hamiltonian dynamics itself.

\paragraph{Restatement: time evolution is a canonical
transformation.}  The map that sends every point to where it will be a
time $t$ later is itself a change of variables on phase space,
$\vec\Gamma(0)\mapsto\vec\Gamma(t)$, with its own Jacobian matrix.
Liouville's theorem says that this map preserves volume, and the
stronger statement is true as well: the map satisfies the symplectic
condition of \S\ref{sec:canonical}, so Hamiltonian time evolution is a
canonical transformation, for every $t$.

This restatement is worth pausing on, because it explains why the
divergence route was taken.  Computing the determinant of the
$2n\times2n$ Jacobian of the time-evolution map directly is hopeless.
Equation \eqref{eq:liouville_theorem} is the infinitesimal version of
the same statement and collapses to two cancelling terms.

\paragraph{What this is used for.}  Two things, in these notes.  In
statistical mechanics, applying the theorem to a probability density
over phase space gives the Liouville \emph{equation}, which is the
starting point of the kinetic-theory chain in \texttt{sidm.tex}
Appendix~A.  In numerical work, it is the property one wants an
integrator to inherit, which is the subject of the next section.

%----------------------------------------------------------------------
\section{Symplectic integrators}
\label{sec:symplectic}

\paragraph{Why the accurate method is the wrong method.}  Suppose we
integrate an orbit numerically.  The obvious choice is a
general-purpose scheme such as fourth-order Runge--Kutta, whose error
after one step of size $h$ is of order $h^5$ --- far smaller than the
$h^3$ of the crude-looking method below.  Run both for a few orbits and
Runge--Kutta wins.  Run both for $10^6$ orbits, which is what a
cosmological simulation does, and Runge--Kutta has lost the answer
entirely: its energy drifts steadily in one direction, and the orbit
spirals in or out.

The reason is that Runge--Kutta is not a canonical map.  Each step
changes phase-space volume slightly, and those changes accumulate in
the same direction.  Local accuracy does not help, because the defect
is structural rather than a matter of step size.

\paragraph{The leapfrog step.}  Split the Hamiltonian into a part
depending only on momenta and a part depending only on coordinates,
which is possible for any system of the form
$\Ham = \sum_i p_i^2/2m_i + V(q)$:
%
\begin{equation}
  \Ham = T(p) + V(q) .
\end{equation}
%
Each part alone can be integrated exactly.  With $\Ham = T(p)$ the
momenta are constant and the coordinates move uniformly; this is called
a \textbf{drift}.  With $\Ham = V(q)$ the coordinates are constant and
the momenta change at a fixed rate; this is a \textbf{kick}.  Both are
exact solutions of a Hamiltonian system, so both are canonical maps.

The drift-kick-drift form of leapfrog composes them:
%
\begin{align}
  q_{1/2} &= q_0 + \frac{h}{2}\,\frac{p_0}{m}
    &&\text{drift for } h/2, \label{eq:dkd_1}\\
  p_1 &= p_0 - h\,\nabla V(q_{1/2})
    &&\text{kick for } h, \label{eq:dkd_2}\\
  q_1 &= q_{1/2} + \frac{h}{2}\,\frac{p_1}{m}
    &&\text{drift for } h/2. \label{eq:dkd_3}
\end{align}
%
The kick-drift-kick form interleaves them the other way and behaves
equivalently; which one a code uses is usually decided by where in the
step the output is written.

\paragraph{Why it works.}  A composition of canonical transformations
is canonical, because the defining brackets \eqref{eq:canonical_test_pb}
are preserved by each factor in turn.  Every leapfrog step is therefore
a canonical map, and phase-space volume is preserved exactly --- not to
order $h^2$, but to machine rounding.  The scheme inherits the property
that \S\ref{sec:liouville} proved for the true dynamics.

The practical consequence follows from a deeper result, known as
backward error analysis: a symplectic integrator applied to $\Ham$ does
not approximately solve $\Ham$, it \emph{exactly} solves a nearby
``shadow'' Hamiltonian $\tilde\Ham = \Ham + \mathcal O(h^2)$.  Since it
conserves $\tilde\Ham$ exactly, and $\tilde\Ham$ differs from $\Ham$ by
a bounded amount, the error in $\Ham$ oscillates within a fixed band
forever instead of growing.  That is the behaviour a long simulation
needs, and it is why every $N$-body code in cosmology integrates this
way.

\paragraph{Two further properties.}  Leapfrog is
\emph{time-reversible}: running the three substeps backwards returns
the starting state exactly, which no Runge--Kutta method does.  It is
also second-order accurate despite requiring only one force evaluation
per step, because the half-step offset between coordinates and momenta
makes the truncation error symmetric, cancelling the first-order term.
The force evaluation is where essentially all the cost of an $N$-body
step lies, so one evaluation per step matters.

\paragraph{Where the assumptions fail.}  Two situations break the
argument, and both appear in practice.  If the timestep is changed
during the run --- as adaptive schemes do --- then the composition is
no longer a fixed canonical map, and the shadow Hamiltonian changes
whenever $h$ changes, which reintroduces drift.  Codes limit the damage
by changing $h$ rarely, in powers of two, and symmetrically.  If the
force is velocity-dependent, as the Coriolis force in a rotating frame
is, the clean split into $T(p) + V(q)$ fails and the kick has to be
handled differently; the standard fix is a Boris rotation placed
symmetrically between the two half kicks, which preserves the
time-centring.

%----------------------------------------------------------------------
\section{Where to read more}
\label{sec:reading}

Goldstein, Poole \& Safko, \emph{Classical Mechanics} (3rd edition), is
the standard graduate text and covers Chapters~2, 8 and 9 of the
material above at length, including the four generating functions and
Hamilton--Jacobi theory.  Landau \& Lifshitz, \emph{Mechanics}, is much
shorter and derives everything from symmetry, which is the more elegant
route once the mechanics is familiar.  Arnold, \emph{Mathematical
Methods of Classical Mechanics}, is the geometric treatment, where the
symplectic condition of \S\ref{sec:canonical} becomes the natural
starting point rather than a technical test.  For the numerical
material of \S\ref{sec:symplectic}, Hairer, Lubich \& Wanner,
\emph{Geometric Numerical Integration}, is the reference, and
Springel's 2005 GADGET-2 paper describes how these ideas are applied in
a cosmological $N$-body code.

Within this collection, \texttt{sidm.tex} Appendix~A continues directly
from \S\ref{sec:liouville} into the Liouville equation, the BBGKY
hierarchy and the Boltzmann equation, and \texttt{qft.tex}~\S3 and
\S7 repeat \S\ref{sec:lagrange} and \S\ref{sec:legendre} for fields
rather than particles.

\end{document}
